
\chapter{Discussion}
\label{ch:discussion}

This chapter reads the measurements of Chapters~\ref{ch:memory_evaluation} to~\ref{ch:spe} as
answers to the three research questions, then states what those answers do not cover. Phone and
workstation results stay separate throughout: the phone gives accuracy and device cost on one
handset, and the workstation tests the same memory modules on more questions, readers and
benchmarks, including the graph repairs and SPE.

Tables~\ref{tab:ch9_claims} and~\ref{tab:ch9_claims_b} list the main claims of
Chapters~\ref{ch:memory_evaluation} to~\ref{ch:spe}, each with its sample size, its test and its
status under the labels defined in Section~\ref{sec:validity}.

\begin{table}[htbp]
  \centering
  \footnotesize
  \renewcommand{\arraystretch}{1.15}
  \begin{tabularx}{\textwidth}{@{}>{\raggedright\arraybackslash}X
      >{\raggedright\arraybackslash}p{0.16\textwidth}
      >{\raggedright\arraybackslash}p{0.15\textwidth}
      >{\raggedright\arraybackslash}p{0.12\textwidth}l@{}}
    \toprule
    Claim & Sample & Test & Status & Section \\
    \midrule
    \multicolumn{5}{@{}l}{\textit{Measurement}} \\
    The automatic judge agrees with a blind human annotator
      & 100 judged answers, 86 questions; 83 decisive labels & Agreement 82/83; Cohen's $\kappa=0.976$
      & Descriptive & \ref{sec:judge} \\
    \midrule
    \multicolumn{5}{@{}l}{\textit{Answer quality on the workstation (RQ1)}} \\
    Vector retrieval beats truncated full context on LoCoMo
      & 200 paired; Gemma E2B, E4B, Qwen 4B & McNemar, Holm; cluster bootstrap
      & Primary & \ref{sec:ch4_vr_fc} \\
    Vector retrieval beats triple-only graph memory (GraphRAG v3)
      & 200 paired; Gemma E2B, E4B & McNemar, Holm & Primary & \ref{sec:ch4_vr_graph} \\
    GraphRAG v3 $\to$ v4 (source turns restored) raises E2B from 21.5\% to 33.0\%
      & 200 paired & McNemar, Holm & Primary & \ref{sec:ch4_source_restoration} \\
    E2B to E4B reader under vector retrieval: +4.5 points, not separated
      & 200 paired & McNemar, Holm & Primary & \ref{sec:ch4_scaling} \\
    Vector retrieval beats truncated full context on LongMemEval-S
      & 60 paired; Gemma E2B & McNemar, Holm & Primary & \ref{sec:ch4_lme} \\
    The same ordering holds for the other readers and the LongMemEval-S graph cells
      & 60--200 paired per reader & McNemar, unadjusted & Exploratory & \ref{sec:ch4_lme} \\
    Long-window full context matches or exceeds vector retrieval on LoCoMo; separated for E2B only
      & 200 paired; Gemma E2B, E4B, Qwen 4B & McNemar, unadjusted & Exploratory
      & \ref{sec:ch4_long_fc} \\
    Vector retrieval stays 5--8 points ahead of long-window full context on LongMemEval-S, not separated
      & 60 paired; same three readers & McNemar, unadjusted & Exploratory
      & \ref{sec:ch4_long_fc} \\
    Retrieval depth $k=2\to50$ adds 15--27 points, more than the tested reader doublings
      & 100 paired; five readers & McNemar, unadjusted; magnitudes not tested
      & Exploratory & \ref{sec:ch4_retrieval_depth} \\
    Question category is associated with correctness
      & Eight Gemma cells of 200 & Permutation $\chi^2$ within conversations, Holm over eight
      & Exploratory & \ref{sec:ch4_locomo_categories} \\
    Reasoning mode gives variable gains; 13--35\% of retrieval-condition outputs unfinished
      & 200 (LoCoMo), 60 (LongMemEval-S) & McNemar, unadjusted & Exploratory
      & \ref{sec:ch4_reasoning} \\
    The official LoCoMo prompt lowers accuracy without changing the ordering
      & 200 paired; Gemma E2B, E4B & McNemar, unadjusted & Exploratory
      & \ref{sec:ch4_robustness} \\
    \bottomrule
  \end{tabularx}
  \caption{Main claims with their evidence: measurement and workstation answer quality (RQ1).
  \emph{Pre-registered}: the criterion was fixed before the run that tests it; only these rows are
  confirmatory. \emph{Primary}: one of the nine paired comparisons Holm-corrected as a family that
  was designated during the analysis. \emph{Exploratory}: a test reported with its unadjusted
  $p$-value. \emph{Descriptive}: a measurement reported without a test. McNemar tests are exact and
  two-sided.}
  \label{tab:ch9_claims}
\end{table}

\begin{table}[htbp]
  \centering
  \footnotesize
  \renewcommand{\arraystretch}{1.15}
  \begin{tabularx}{\textwidth}{@{}>{\raggedright\arraybackslash}X
      >{\raggedright\arraybackslash}p{0.16\textwidth}
      >{\raggedright\arraybackslash}p{0.15\textwidth}
      >{\raggedright\arraybackslash}p{0.12\textwidth}l@{}}
    \toprule
    Claim & Sample & Test & Status & Section \\
    \midrule
    \multicolumn{5}{@{}l}{\textit{Phone resources (RQ2)}} \\
    Embedding prefixes have no detectable effect on phone vector accuracy (E2B 10$\to$8, E4B 8$\to$10)
      & 20 paired; Gemma E2B, E4B & McNemar, unadjusted & Exploratory & \ref{sec:retrieval} \\
    Workstation accuracy cannot stand in for phone accuracy (83\% agreement, threshold 90\%)
      & 100 paired; Gemma E2B vector & Agreement threshold; $\kappa$; McNemar
      & Pre-registered & \ref{sec:ch5_parity} \\
    Vector retrieval has higher accuracy, lower energy per answer and shorter query time than full
      context in all five readers
      & 20 questions per cell; 3 vector sittings per reader; 1--3 full-context sittings & None
      & Descriptive & \ref{sec:ch5_phone_comparison} \\
    5.8--26.9$\times$ more correct answers per kilojoule
      & As above; ratios within matched sittings & Resampling sensitivity: lower end 2.6--8.8$\times$
      & Descriptive & \ref{sec:ch5_charge} \\
    Peak RAM and temperature do not follow the energy ordering
      & As above & None & Descriptive & \ref{sec:ch5_resources} \\
    The early phone graph build (256-token limit, 58/84 batches failed) cost about 26$\times$ more
      energy per turn than vector ingestion
      & One graph build against three vector sittings & None; range 22--28$\times$ by sitting
      & Descriptive & \ref{sec:ch5_graph_boundary} \\
    \midrule
    \multicolumn{5}{@{}l}{\textit{Graph memory and SPE (RQ3)}} \\
    Small Qwen extractors fail on 79--85\% of batches; more retrieved facts recover 1--2 answers
      & 1,179 batches; 200 paired & McNemar, unadjusted & Exploratory & \ref{sec:ch6_extractors} \\
    The graph producer, more than the reader, sets GraphRAG v3 accuracy
      & Four cells of 200 & None & Descriptive & \ref{sec:ch6_extractors} \\
    SPE misses Criterion~1 and stays below GraphRAG v4
      & 200 paired; E2B (E4B addendum) & Registered thresholds; McNemar, Holm vs v4
      & Pre-registered & \ref{sec:spe_results} \\
    SPE passes Criterion~2 numerically, without a gain separable from GraphRAG v3
      & 200 paired; E2B & Registered thresholds; McNemar vs v3
      & Pre-registered & \ref{sec:spe_results} \\
    SPE-2 selects no source turn
      & 800 prompts & Count & Descriptive & \ref{sec:spe_why_triples} \\
    A per-question choice between GraphRAG v3 and v4 has headroom (oracle)
      & 200; uses judged labels & None; post hoc bound & Descriptive & \ref{sec:spe_discussion} \\
    \bottomrule
  \end{tabularx}
  \caption{Main claims with their evidence: phone resources (RQ2) and graph memory and SPE (RQ3).
  Status labels as in Table~\ref{tab:ch9_claims}. Phone rows use phone-measured accuracy only.}
  \label{tab:ch9_claims_b}
\end{table}

\section{RQ1: how memory strategy affects answer quality}
\label{sec:ch9_rq1}

Under a 4,096-token history budget, the strategy ordering was stable across every reader tested. On
the phone, vector retrieval
answered 20\% to 50\% of the 20 LoCoMo questions correctly, against 0\% to 10\% for full context.
On the workstation's 200-question LoCoMo set it had the highest accuracy point estimate for all
five readers, 28.5 to 32.5 points above truncated full context in the primary comparisons.
Triple-only graph memory sat between the two, trailing vector retrieval by 13.0 to 15.5 points
where extraction worked. LongMemEval-S repeated the ordering: vector retrieval scored 26.7\% to
56.7\%, no truncated full-context cell exceeded 10\%, and the repaired Gemma graphs again sat between the two
(Section~\ref{sec:ch4_lme}).

For a system designer, the sizes of the two settings behind that ordering are also informative.
Raising retrieval depth from two to fifty turns raised complete-evidence coverage from 16\% to 64\%
and accuracy by 15 to 27 points in every reader. The tested reader doublings were worth 4.5 points
from Gemma~4~E2B to E4B and 6.5 points from Qwen~3.5~2B to 4B. For readers of 0.8 to 4 billion
parameters under a 4,096-token history budget, what the memory system put in front of the reader
therefore moved accuracy more than the tested increases in reader size. The comparison sets a
25-fold range of depth against one doubling of the reader, and the larger depths cost longer
prompts. It does not rank the two settings in general, and it does not apply to readers that can
hold the entire history.

Two qualifications belong with this answer. First, full context scores low because the history
does not fit: the workstation truncated LoCoMo histories to a 4,096-token window, and the phone kept
only the most recent turns (79 of 419 in the conversation of
Figure~\ref{fig:ch3_benchmark_anatomy}). With a long workstation window, full context rose to
41.5--44.0\% for the three readers tested on LoCoMo, where it held the whole history, at or above
vector retrieval, and to 40.0--48.3\% on LongMemEval-S, where vector retrieval kept a 5- to 8-point
lead that 60 questions could not separate (Section~\ref{sec:ch4_long_fc}). Those long-window
prompts were about 39 and 48 times longer than the retrieval prompts. The vector advantage is
therefore an advantage in evidence per prompt token under a small context budget; on LoCoMo, a
reader with a long enough window matched or passed it. On the phone that window is the constraint: a rough extrapolation from
the measured query time per prompt token puts a whole LoCoMo history at about 45 minutes per
question, a length never run there. Second, the ordering holds for overall
accuracy, not in every category: temporal questions never exceeded 26\% in the main LoCoMo grid,
and none of the tested memory designs reliably repaired that weakness.

\section{RQ2: what memory strategy costs on a phone}
\label{sec:ch9_rq2}

On the Galaxy S23 Ultra, the quality and cost orderings agreed. Vector retrieval had a higher
accuracy point estimate, lower energy per answer and shorter query time than truncated full
context in all five reader pairs, and 5.8 to 26.9 times more correct answers per kilojoule where
that ratio is finite. On this slice, every question that full context answered correctly was also answered correctly by
vector retrieval. Bounded retrieval was therefore not a compromise between quality and cost on this
device: it was better on both.

That result is whole-workload efficiency: each vector sitting's energy includes building its
memory, about 60\% of an E2B vector sitting. Construction is also where the strategies separate
most sharply. In the phone pilot, extracting a graph cost about 26 times more energy per dialogue
turn than embedding (Table~\ref{tab:ch5_construction}). This cost is paid before the first question
and is invisible to prompt-token accounting. Server-side evaluations of memory frameworks likewise find ingestion and
retrieval dominating processing time~\citep{ai_hippocampus}. Device comparisons should therefore price construction
separately from each query, as the ledger in Chapter~\ref{ch:on_device_memory} does.

The resource results do not collapse into one ordering. Peak system RAM and peak hottest-zone
temperature did not follow energy and latency, so a policy that optimizes only joules can still
violate a memory or thermal limit. This is why the individual traces are reported alongside the
joint measure.

Two boundaries limit how far these results travel. The workstation-to-phone parity gate failed at
83\% agreement against a 90\% threshold, so every phone claim rests on the 20-question phone
cells. The graph strategy appears on the phone only as one early construction pilot, with 58 of 84
extraction batches failing, so this thesis makes no claim about the accuracy of a repaired
on-device graph.

\section{RQ3: what graph memory can and cannot do here}
\label{sec:ch9_rq3}

Here the mechanism, rather than the ranking, is the contribution. The repair sequence separated
four places where evidence is lost: facts that never entered the graph because extraction failed or
omitted them, stored facts excluded by the retrieval rule or triple budget, selected triples missing
detail from their source turn, and evidence in the prompt that the reader did not use. Restoring the
source text behind the selected triples raised Gemma~4~E2B from 21.5\% to 33.0\% ($p=0.000170$ after
the thesis's Holm correction) and E4B from 27.5\% to 32.5\%.

Both conditions started from the same ten retrieved turns, so the E2B gain points to the triple
representation rather than retrieval as the main place the loss arises. For E4B, restoration
recovered 5 of the 13 points that separated triples from vector retrieval.

SEEM and LeanMem also return from a compact memory unit to its source text
(Section~\ref{sec:ch2_graphs}): SEEM for every retrieved event frame~\citep{seem}, and LeanMem when a
planner judges that a question needs detail~\citep{leanmem}. Neither reports energy. This thesis
measures how much that return is worth for small readers under a small budget, and the device cost of
building the compact unit in the first place.

SPE tried to keep the restoration gain while controlling its token cost. Criterion~1 failed, and
Criterion~2 cleared its numerical thresholds without a distinguishable improvement over version~3;
none of the 800 pre-specified SPE-2 prompts selected a source turn. Under per-token scoring, an
upgrade to source text had to earn roughly four times what a competing triple offered, and the
budget ran out on triples first (Section~\ref{sec:spe_why_triples}). The utilities were not logged,
so this explanation follows from the rule and the measured costs rather than from a logged trace. A post hoc
oracle choosing between v3 and v4 per question with judged labels reaches 35.0\% for E2B and 33.5\%
for E4B at about 560 total prompt tokens. It is an upper bound, not an implementable policy, but it
shows headroom at the question level, so the negative result applies to the tested per-fact rule
rather than to selective expansion in general.

The extraction floor sets the outer limit. The smaller Qwen extractors failed on roughly four batches
in five and the phone pilot on 58 of 84, so for those configurations the question is less which
retrieval rule to use than how much of the history extraction captures. Where extraction worked,
raising the graph retrieval budget from twelve to twenty-four facts was worth 2.0 to 8.0 points;
where it had failed, the same change moved one or two questions out of 200.

\section{What it took to make these measurements trustworthy}
\label{sec:ch9_method}

Three lessons from the measurement procedure apply beyond this study.

\paragraph{Device accuracy has to be measured on the device.}
The pre-registered parity check failed, so phone accuracy comes only from phone answers and is never
substituted from the workstation. The smaller phone cells are the price of that separation.

\paragraph{A budgeted selection rule needs a defined tie-break.}
The first graph retrieval rule gave most edges of a matched speaker hub the same score, so the facts
that filled the budget depended on the process's hash seed; 164 of 200 version-1 questions had such a
tie at the cutoff. The vector-first retrieval used from version~2 onward returned identical contexts
under different seeds, so the record results are unaffected.

\paragraph{The energy measurement is simple and cross-checked, with one known gap.}
Energy comes from the phone's charge counter, weighted by the measured voltage. It agrees with the
reported battery-level drop, and with integrated current to within 10\% on sittings longer than 50
minutes. These are consistency checks, not independent instruments. Sittings that started at the
highest possible charge (100\%) under-count because the counter pauses at full charge; Appendix~\ref{app:on_device_statistics}
quantifies the effect, and no comparison changes direction.

\section{Threats to validity and what the results do not show}
\label{sec:ch9_threats}
\label{sec:ch9_not_shown}

The chapter-level threat sections state the specific limits; the consolidated ones are these.

\paragraph{Sample and independence.}
The 200 LoCoMo questions come from ten conversations and the 60 LongMemEval-S questions from a
stratified sample of one benchmark, so question-level observations are not independent. Paired tests
and cluster sensitivity checks preserve the direction of the main comparisons, and category cells of
50 or 10 questions support only coarse statements.

\paragraph{Judging.}
All reported accuracies come from one automatic judge. Re-judging identical predictions changed no
verdict within one judging session and one verdict in 100 across sessions. A blind human audit of
100 judged answers agreed with the judge on 82 of the 83 decisive labels (Cohen's $\kappa=0.976$),
with no correct answer marked wrong; 17 answers were too ambiguous to label, and counting them all as
disagreements would lower agreement to 82 of 100 (Section~\ref{sec:judge}). The audit had one
annotator, so it does not measure agreement between humans, and no second judge model was compared.
The primary strategy gaps are far larger than the variation from re-judging. The audit counts only
decisive labels from a sample drawn across conditions, so it does not bound judge bias within any
single strategy comparison.

\paragraph{One device, one runtime, uneven replication.}
All device measurements come from one Galaxy S23 Ultra, one Android build and a CPU-oriented
llama.cpp configuration, with three sittings for each vector cell and complete full-context
replication only for Gemma~4~E2B. Each phone cell has 20 questions, so the efficiency multipliers
are uncertain: resampling the question pairs puts their lower end at 2.6 to 8.8$\times$
(Section~\ref{sec:ch5_charge}). A different SoC, an accelerator path or a different thread and
governor configuration may change the ordering.

\paragraph{Cross-study comparison.}
External accuracy numbers in this thesis are context, not comparators. Readers, judges, retrieval
units, context limits and question populations differ across the cited studies, so their effect sizes
are not pooled with those measured here.

Four statements are excluded deliberately, because the evidence in this thesis does not support them.

\begin{enumerate}
  \item \textbf{That knowledge graphs fail for conversational memory.} The evidence concerns graphs
        extracted by small local models under a small generation budget and queried by turn-anchored
        triple selection, without multi-hop traversal; in these speaker-centred graphs most edges already
        lie next to the retrieved facts' entities. Community-summary designs, much larger extractors
        and curated knowledge graphs were not evaluated. SmartRAG demonstrates a graph-based mobile
        memory system built with a distilled extractor~\citep{smartrag}, and the token-matched
        comparison of \citet{rag_vs_graphrag} retains a temporal-accuracy advantage for graph
        retrieval at server scale.
  \item \textbf{That selective expansion cannot work.} The tested per-fact rule failed, and the
        oracle shows that a question-level choice had headroom on the same cells. A query-level
        policy of the kind LeanMem applies~\citep{leanmem} was not tested here.
  \item \textbf{That the device draws a fixed power.} The pilot's average of about 4~W is the
        quotient of one workload's energy and its duration. The frequency and temperature traces are
        consistent with thermal or power throttling; they do not identify a governor policy or
        establish an invariant power figure.
  \item \textbf{That prefill dominates the measured query time.} The phone logs record total query
        time, and prefill and decode were not timed separately. The relationship between prompt
        length and query time is consistent with prompt processing dominating, which is weaker than
        a measured phase split.
\end{enumerate}

Stated positively, for conversational question answering with small local readers under a small
context budget, retrieving a bounded amount of source text was the strongest quality-cost point
measured, and on the workstation it stayed within 1.5 to 8 points of reading the whole history on
LoCoMo, and ahead of long-window full context on LongMemEval-S, with 2 to 3\% of the prompt tokens; the compact graph representation lost evidence that its own source spans still contained;
restoring that text recovered part of the loss at workstation scale; the tested rule for spending a
token budget on that restoration did not work; and building the compact representation carries a
device cost that prompt-level accounting does not show.
